\documentclass[draft]{cpkkuproject}

\titleTH{การพัฒนาระบบแปลงเสียงเป็นข้อความโดยใช้ปัญญาประดิษฐ์เพื่อบันทึกข้อมูลอัตโนมัติในงานพยาธิวิทยา}
\titleEN{Development of an AI-Based Speech-to-Text System for Automatic Data Recording in Pathology}

\firstStudentID{663380271-2}
\firstStudentPrefixTH{นาย}
\firstStudentPrefixEN{Mr.}
\firstStudentNameTH{ธนินท์รัฐ โลหะสาร}
\firstStudentNameEN{Thaninrat Lohasan}

\secondStudentID{663380304-3}
\secondStudentPrefixTH{นาย}
\secondStudentPrefixEN{Mr.}
\secondStudentNameTH{เจษฎา กลางถิ่น}
\secondStudentNameEN{Chetsada Klangthin}

\degreeTH{วิทยาศาสตรบัณฑิต}
\degreeEN{Bachelor of Science}
\majorTH{ปัญญาประดิษฐ์}
\majorEN{Artificial Intelligence}
\advisorTH{ผศ. ดร.ไอศูรย์ กาญจนสุรัตน์}
\coadvisorTH {ผศ. นพ.ชัยวัชร์ อภิวาทนสิริ}
\advisorEN {Asst. Prof. Isoon Kanjanasurat} 
\coadvisorEN {Asst. Prof. Chaiwat Apiwatnasiri}

\academicYearTH{2569}
\academicYearEN{2026}

\keywordsTH{ปัญญาประดิษฐ์, การแปลงเสียงเป็นข้อความ, พยาธิวิทยา, Whisper, การสกัดข้อมูลทางการแพทย์, การกรองสัญญาณรบกวน}
\keywordsEN{Artificial Intelligence, Speech-to-Text, Gross Pathology, Whisper, Clinical Information Extraction, Noise Reduction}

\usepackage{longtable,array}
\usepackage{hyperref}
\hypersetup{
    colorlinks=true,
    citecolor=black,  % เปลี่ยนสีตัวอ้างอิงในเนื้อหา (เช่น [3], [6]) ให้เป็นสีดำ
    linkcolor=black,  % สีของลิงก์ภายในเอกสาร (สารบัญ, สมการ) เป็นสีดำ
    urlcolor=black    % สีของลิงก์เว็บไซต์เป็นสีดำ
}
\begin{document}
\maketitlepage 
\pagenumbering{arabic}
\setcounter{page}{0}
\renewcommand{\thepage}{\ThaiAlph{\value{page}}}

\begin{abstractTH}
โครงงานนี้นำเสนอกระบวนการพัฒนาระบบรู้จำเสียงพูดและสกัดข้อมูลทางการแพทย์แบบออฟไลน์สำหรับงานตรวจชิ้นเนื้อมะเร็งเต้านมบนคอมพิวเตอร์ส่วนบุคคล เพื่อแปลงเสียงบรรยายการตรวจชิ้นเนื้อเป็นข้อความและสกัดข้อมูลลงแบบฟอร์มมาตรฐาน 15 ฟิลด์ตามเกณฑ์ CAP โดยอัตโนมัติ ช่วยลดอุปสรรคจากการสวมถุงมือเปื้อนสารเคมี ลดข้อผิดพลาดจากการบันทึกข้อมูล และคุ้มครองความเป็นส่วนตัวตามมาตรฐาน PDPA และ HIPAA ระบบที่พัฒนาขึ้นประกอบด้วย 4 ส่วนหลัก ได้แก่ (1) กลไกคัดกรองเสียงรบกวนอัตโนมัติ (Blind SNR Gate) (2) การกำหนดข้อความนำบริบทเชิงเปรียบเทียบเสียงพ้อง ($P_4$) (3) การปรับแต่งแบบจำลอง Whisper Small ด้วยเทคนิค LoRA ร่วมกับการลดทอนความละเอียดเป็น 8 บิต บน CPU และ (4) ตัวสกัดข้อมูลลงแบบฟอร์มมาตรฐานร่วมกับกลไกแจ้งเตือนความปลอดภัยหลายสัญญาณ (Multi-Signal Safety Flagging) ผลการทดสอบบนชุดข้อมูลจำลองสภาวะทางคลินิกจำนวน 1,000 กรณีศึกษา พบว่าระบบสามารถคัดแยกสัญญาณเสียงรบกวนได้ถูกต้องร้อยละ 98.67 และลดอัตราความผิดพลาดวิกฤตระดับเคส ($\text{CER}_{\text{case}}$) บนชุดทดสอบสำนวนใหม่ลงเหลือร้อยละ 5.33 พร้อมทำคะแนนเอฟวันเฉลี่ยทั้ง 15 ฟิลด์ได้ร้อยละ 99.13 ส่วนชุดทดสอบผู้พูดและคำศัพท์ใหม่สามารถลด $\text{CER}_{\text{case}}$ ลงเหลือร้อยละ 11.33 โดยใช้เวลาประมวลผลเฉลี่ยเพียง 6.07--6.57 วินาทีต่อกรณี ควบคู่กับกลไกแจ้งเตือนความปลอดภัยที่จำกัดข้อผิดพลาดหลุดรอดได้อย่างมีประสิทธิภาพ โครงงานนี้จึงเป็นต้นแบบเครื่องมือสนับสนุนการจัดทำร่างรายงานการตรวจชิ้นเนื้อเบื้องต้นสำหรับพยาธิแพทย์ที่ทำงานแบบออฟไลน์ได้อย่างถูกต้อง รวดเร็ว ทนทานต่อเสียงรบกวน และมีความปลอดภัยสูงทางคลินิก
\end{abstractTH}


\begin{abstractEN}
Gross examination in surgical pathology requires pathologists to dictate complex anatomical measurements while wearing specimen-contaminated gloves in an acoustic environment degraded by formalin fume-hood exhaust noise, rapid speech, and mid-sentence self-corrections. This project presents PathoWhisper, an offline clinical speech recognition and structured information extraction system for breast cancer gross pathology running entirely on local personal computer CPUs. To address fundamental challenges in acoustic robustness, prompt design, domain adaptation, and clinical safety, this work introduces four core technical contributions: (1) a Blind SNR-Gated Adaptive Denoising module ($\widehat{\mathrm{SNR}} = P_{85} - P_{15}, \tau^* = 13.55\text{ dB}$) paired with a calibrated FFT spectral subtraction filter (\texttt{afftdn}, $\mathrm{nr} = 6\text{ dB}$) that prevents consonant over-suppression on clean speech while restoring intelligibility under severe fume-hood noise; (2) a systematic Prompt Structure and Token-Budget Ablation ($P_0 \dots P_4$) demonstrating orthogonal decomposition between clinical terminology errors and numerical dimension formatting; (3) Parameter-Efficient Fine-Tuning of Whisper Small via Low-Rank Adaptation (LoRA, $r=8, \alpha=32$ across 48 decoder attention projections, training only 589,824 parameters or 0.2434\% of the model) merged and quantized into CTranslate2 INT8 for real-time CPU execution; and (4) a deterministic 15-field CAP extractor coupled with Multi-Signal Clinical Safety Flagging.

Evaluation was conducted under a strict Zero-Phrase-Leakage protocol across a 1,000-case dataset partitioned into training ($N=600$), validation ($N=100$), seen-speaker unseen-phrasing test (\texttt{test\_seen}, $N=150$), and held-out speaker/vocabulary test (\texttt{test\_unseen}, $N=150$) splits covering all 15 standardized CAP breast cancer fields. The blind SNR gate achieved 98.67\% routing accuracy (296/300 test cases) and significantly reduced case-level critical error rate ($\mathrm{CER}_{\mathrm{case}}$) compared to static always-on \texttt{afftdn} filtering from 18.00\% to 13.00\% ($p = 0.0015$). Three-epoch LoRA fine-tuning converged smoothly ($\mathrm{train\_loss}: 0.4249 \rightarrow 0.1318$, $\mathrm{val\_loss}: 0.2747 \rightarrow 0.2189$). Compared to the zero-shot baseline ($S_0$), the full PathoWhisper pipeline ($S_3$) reduced $\mathrm{CER}_{\mathrm{case}}$ on \texttt{test\_seen} from 12.67\% (19/150) to 5.33\% (8/150, McNemar $p = 0.0127$) with a 15-field Macro-$F_1$ of 99.13\%, and more than halved $\mathrm{CER}_{\mathrm{case}}$ on \texttt{test\_unseen} from 24.00\% (36/150) to 11.33\% (17/150, McNemar $p = 6.60 \times 10^{-5}$), while reducing Clinical Term Error Rate (C-TER) from 12.03\% to 4.05\% and Numerical Dimension Error Rate (N-DER) from 5.98\% to 1.58\%. Furthermore, the Multi-Signal Clinical Safety Flagging mechanism achieves up to 100.0\% error recall on zero-shot baselines and constrains unflagged silent critical case error rates to 1.83\%--1.96\% on \texttt{test\_seen} and 5.00\%--7.77\% on \texttt{test\_unseen} under LoRA adaptation, while reducing physician review workload to 27.33\%--33.33\% at an average CPU latency of 5.75--6.57 seconds per case ($\mathrm{RTF} \approx 0.23\text{--}0.25$).
\end{abstractEN}


\begin{acknowledgement}
ผู้จัดทำโครงงานขอขอบพระคุณ ผู้ช่วยศาสตราจารย์ ดร.ไอศูรย์ กาญจนสุรัตน์ อาจารย์ที่ปรึกษาโครงงานหลัก ที่ได้กรุณาสละเวลาให้คำแนะนำ คำปรึกษา ถ่ายทอดองค์ความรู้ และชี้แนะแนวทางการดำเนินงานวิจัยอันเป็นประโยชน์อย่างยิ่งตลอดระยะเวลาการดำเนินโครงงาน ตลอดจนช่วยตรวจแก้ไขรายงานโครงงานฉบับนี้ให้มีความสมบูรณ์ถูกต้อง

ขอขอบพระคุณ ผู้ช่วยศาสตราจารย์ นายแพทย์ชัยวัชร์ อภิวาทนสิริ อาจารย์ที่ปรึกษาร่วม สาขาวิชาพยาธิวิทยา คณะแพทยศาสตร์ มหาวิทยาลัยขอนแก่น ที่ได้กรุณาให้คำปรึกษาและข้อเสนอแนะเชิงลึกด้านขั้นตอนการตรวจชิ้นเนื้อทางศัลยพยาธิวิทยาด้วยตาเปล่า (Gross Examination) คำศัพท์เฉพาะทางพยาธิวิทยา ตลอดจนมอบแบบฟอร์มมาตรฐานตรวจสิ่งส่งตรวจเต้านมประจำสาขาวิชาพยาธิวิทยา ซึ่งพัฒนาขึ้นตามเกณฑ์สากลของ College of American Pathologists (CAP) เพื่อนำมาเป็นต้นแบบในการพัฒนาแบบฟอร์มดิจิทัล 15 ฟิลด์ของโครงงานนี้

ขอขอบพระคุณคณาจารย์ประจำหลักสูตรปัญญาประดิษฐ์และวิทยาลัยการคอมพิวเตอร์ มหาวิทยาลัยขอนแก่น ทุกท่าน ที่ได้ประสิทธิ์ประสาทวิชาความรู้และทักษะทางวิชาการอันเป็นรากฐานสำคัญในการศึกษาค้นคว้า ตลอดจนขอขอบคุณเพื่อนร่วมชั้นเรียนและผู้มีส่วนเกี่ยวข้องทุกท่านที่ได้ให้ความช่วยเหลือ ร่วมแลกเปลี่ยนความคิดเห็น และเป็นกำลังใจในการดำเนินงานร่วมกันมาโดยตลอด

ท้ายที่สุดนี้ ผู้จัดทำขอกราบขอบพระคุณบิดา มารดา และสมาชิกในครอบครัวทุกท่าน ผู้ซึ่งเป็นแรงผลักดันและคอยให้การสนับสนุนในทุกด้านด้วยความรักและความเข้าใจเสมอมา จนกระทั่งโครงงานฉบับนี้สำเร็จลุล่วงได้ด้วยดี


\end{acknowledgement}

\tableofcontents
\clearpage
\listoffigures
\clearpage
\listoftables
\clearpage
\pagenumbering{arabic}
\setcounter{page}{1}

\chapterone
\chaptertwo
\chapterthree
\chapterfour
\chapterfive

\providecommand{\todo}[1]{\textbf{[TODO: #1]}}
% ============================
% ภาคผนวก
% ============================
\setcounter{chapter}{-1}

\appendiceschapter{การตั้งค่าระบบและพารามิเตอร์ของไปป์ไลน์}
\section*{คำสั่งและตรรกะการกรองสัญญาณเสียงแบบปรับตัว (\texttt{D8\_snr\_gated})}

ระบบประเมินค่าอัตราส่วนสัญญาณต่อเสียงรบกวนแบบไม่ต้องใช้สัญญาณอ้างอิง ($\widehat{\mathrm{SNR}}$) จากผลต่างเปอร์เซ็นไทล์ค่ารากกำลังสองเฉลี่ยของพลังงาน (Root-Mean-Square Energy: RMS) ของเฟรมเสียงความยาว 25 มิลลิวินาที (ระยะเลื่อนเฟรม 10 มิลลิวินาที) เปรียบเทียบกับค่าเกณฑ์ขีดแบ่งที่สอบเทียบด้วยดัชนียูเดน (Youden's Index) บนชุดข้อมูลตรวจสอบ ($\tau^* = 13.55\text{ dB}$) โดยจะเรียกใช้งานสายโซ่ตัวกรองความถี่ของโปรแกรม FFmpeg เฉพาะเมื่อ $\widehat{\mathrm{SNR}} < 13.55\text{ dB}$ และส่งผ่านสัญญาณเสียงต้นฉบับโดยไม่กรอง (Bypass) เมื่อ $\widehat{\mathrm{SNR}} \ge 13.55\text{ dB}$ เพื่อรักษาสเปกตรัมความถี่ของเสียงพยัญชนะสั้น

\begin{table}[H]
\centering
\caption{พารามิเตอร์ของกลไกคัดกรองสัญญาณรบกวนแบบปรับตัว (\texttt{D8\_snr\_gated})}
\label{tab:appfilter}
\small
\begin{tabular}{|l|l|p{6.5cm}|}
\hline
\textbf{องค์ประกอบ} & \textbf{พารามิเตอร์} & \textbf{ค่าที่กำหนดและเหตุผลทางวิศวกรรม} \\
\hline
Blind SNR Gate & $\tau^*$ ($P_{85} - P_{15}$ Frame Energy) & $13.55\text{ dB}$ (สอบเทียบจากชุดตรวจสอบ \texttt{val} ให้ค่า $J = 1.0000$, คัดแยกถูกต้องร้อยละ $98.67$ บนชุดทดสอบ) \\
\hline
\texttt{highpass} & ความถี่ตัด ($f_c$) & $80\text{ Hz}$ (กรองตัดเสียงรบกวนความถี่ต่ำจากมอเตอร์พัดลมตู้ดูดควันก่อนเข้าสู่ตัวกรอง \texttt{afftdn}) \\
\hline
\texttt{afftdn} & \texttt{nr} / \texttt{nf} / \texttt{nt} & $6\text{ dB}$ / $-30\text{ dB}$ / \texttt{white} (ระดับการลดเสียงรบกวนแบบนุ่มนวลที่ให้ผลดีกว่าค่าเดิม $12\text{ dB}$ จากการค้นหาแบบกริด) \\
\hline
\texttt{silenceremove} & ต้นไฟล์และระหว่างไฟล์ & ตัดช่วงสัญญาณเงียบที่ต้นไฟล์เกิน $0.1\text{ s}$ และช่วงเงียบระหว่างพูดที่เกิน $0.6\text{ s}$ (ระดับพลังงานต่ำกว่า $-40\text{ dB}$) \\
\hline
\end{tabular}
\end{table}

\section*{ค่าตั้งค่าการปรับแต่งค่าน้ำหนักแบบจำลอง (LoRA) และเอนจินถอดความ}

\begin{table}[H]
\centering
\caption{ค่าตั้งค่าของแบบจำลอง PathoWhisper Small LoRA และเอนจินการถอดความแบบปรับโหมดตามความยาวเสียง}
\label{tab:appconfig}
\small
\begin{tabular}{|l|p{8.2cm}|}
\hline
\textbf{พารามิเตอร์} & \textbf{ค่าที่กำหนด} \\
\hline
แบบจำลองฐานและเอนจินอนุมาน & OpenAI Whisper Small แปลงเข้าสู่เอนจิน CTranslate2 (Faster-Whisper) \\
\hline
การปรับแต่งด้วยพารามิเตอร์จำนวนน้อย (LoRA) & อันดับเมทริกซ์ $r = 8$, สัมประสิทธิ์มาตราส่วน $\alpha = 32$, $\text{dropout} = 0.05$ บนชั้น \texttt{q\_proj, v\_proj} ของตัวถอดรหัส 48 โมดูล (พารามิเตอร์ฝึกฝน $589,824$ ตัว หรือร้อยละ $0.2434$) \\
\hline
ชุดข้อมูลฝึกฝนและตรวจสอบ & \texttt{train} $= 600$ กรณี, \texttt{val} $= 100$ กรณี (ฝึกฝน 3 รอบ, อัตราการเรียนรู้ $\text{lr} = 1 \times 10^{-4}$, ความละเอียดผสม FP16) \\
\hline
อุปกรณ์และชนิดข้อมูลขณะอนุมานผล & หน่วยประมวลผลกลาง CPU (\texttt{cpu\_threads=8}), ลดทอนความละเอียดเป็นเลขจำนวนเต็ม 8 บิต \texttt{INT8} (ขนาดไฟล์ประมาณ $253\text{ MB}$) \\
\hline
พารามิเตอร์การถอดรหัสมาตรฐาน ($\le 27.5\text{ s}$) & \texttt{beam\_size=1}, \texttt{word\_timestamps=True}, \texttt{condition\_on\_previous\_text=False} \\
\hline
การถอดรหัสไฟล์เสียงความยาวสูง ($> 27.5\text{ s}$) & \texttt{without\_timestamps=True}, \texttt{condition\_on\_previous\_text=False}, \texttt{repetition\_penalty=1.15}, \texttt{vad\_filter=True} (\texttt{min\_silence\_duration\_ms=400}) \\
\hline
\end{tabular}
\end{table}

\section*{การออกแบบข้อความนำบริบทแบบแยกส่วนเชิงตั้งฉาก ($P_0 \dots P_4$)}

ข้อความนำบริบทที่ใช้งานในระบบบูรณาการ ($P_4$: \texttt{P4\_phonetic\_contrastive}) ได้รับการออกแบบโดยไม่คัดลอกสำนวนประโยคตัวอย่างจากชุดข้อมูลทดสอบ (Zero-Phrase-Leakage) และมีความยาว 118 โทเคน (อยู่ภายในขีดจำกัด 224 โทเคนของตัวถอดรหัส Whisper) ดังนี้:

\begin{quote}
\small
\texttt{Gross pathology breast specimen (CAP protocol). Terminology \& homophone disambiguation: right breast vs left breast (not write), modified radical mastectomy, simple mastectomy, skin ellipse, nipple everted, inverted, retracted, healed biopsy scar, biopsy cavity with residual mass, infiltrative mass vs no discrete mass (not discreet), entirely fibrocystic change (not fibrousistic), ductal carcinoma in situ with microinvasion, upper outer, upper inner, lower outer, lower inner, central quadrant, deep surgical margin (not lymph or weight margin), superior, inferior, medial, lateral margin, axillary lymph nodes. Format: S-00-0000, 0.0 x 0.0 x 0.0 cm.}
\end{quote}

\appendiceschapter{ฟิลด์ข้อมูลที่สกัดและผลการประเมินโดยละเอียด}
\section*{รายการฟิลด์มาตรฐานทั้ง 15 ฟิลด์}
\begin{table}[H]
\centering
\caption{ฟิลด์ข้อมูลทางศัลยพยาธิวิทยาเต้านมตามเกณฑ์มาตรฐานทั้ง 15 ฟิลด์ที่ระบบรองรับและใช้ประเมินผล}
\label{tab:appfields}
\small
\begin{tabular}{|c|p{4.6cm}|p{7.0cm}|}
\hline
\textbf{ที่} & \textbf{รหัสฟิลด์} & \textbf{ความหมายและบทบาทในการประเมิน} \\
\hline
1 & \texttt{s0\_surgical\_no} & เลขที่สิ่งส่งตรวจทางศัลยกรรม (เช่น \texttt{S-26-8450}) \\
\hline
2 & \texttt{s1\_side} & ด้านของเต้านม (\texttt{Right} / \texttt{Left}) --- \textbf{ฟิลด์วิกฤต (Critical)} \\
\hline
3 & \texttt{s2\_proc} & ประเภทหัตถการผ่าตัด (\texttt{Modified radical mastectomy}, \texttt{Simple mastectomy}) \\
\hline
4 & \texttt{s3\_dims} & ขนาด 3 มิติของสิ่งส่งตรวจหลัก (ซม.) --- \textbf{ฟิลด์วิกฤต (Critical)} \\
\hline
5 & \texttt{s4\_skin} & การมีชิ้นผิวหนังรูปวงรี (\texttt{Skin ellipse}) ติดมากับสิ่งส่งตรวจ \\
\hline
6 & \texttt{s5\_dims} & ขนาด 2 มิติของชิ้นผิวหนัง (ซม.) \\
\hline
7 & \texttt{s6\_nipple} & สภาพหัวนมและลานนม (\texttt{everted}, \texttt{inverted}, \texttt{retracted}) (มีกรณีผลบวกในหมวดหมู่ที่ 4) \\
\hline
8 & \texttt{s7\_biopsy\_scar} & การพบรอยแผลเป็นจากการเจาะตรวจเดิม (\texttt{healed biopsy scar}) (มีกรณีผลบวกในหมวดหมู่ที่ 4) \\
\hline
9 & \texttt{s8\_cavity} & การพบโพรงจากการผ่าตัดครั้งก่อน (\texttt{biopsy cavity}) (มีกรณีผลบวกในหมวดหมู่ที่ 10) \\
\hline
10 & \texttt{s9\_residual\_mass} & ขนาดก้อนเนื้องอกหลงเหลือข้างโพรงผ่าตัดเดิม (มีกรณีผลบวกในหมวดหมู่ที่ 10) \\
\hline
11 & \texttt{s10\_infiltrative} & ลักษณะรอยโรคหลัก (\texttt{Infiltrative}, \texttt{Circumscribed}, \texttt{No mass}) --- \textbf{ฟิลด์วิกฤต (Critical)} \\
\hline
12 & \texttt{s10\_inf\_dims} & ขนาด 3 มิติของก้อนเนื้องอกหลัก (ซม.) --- \textbf{ฟิลด์วิกฤต (Critical)} \\
\hline
13 & \texttt{s10\_5\_quadrant} & ตำแหน่งจตุภาคที่พบก้อนเนื้องอก (\texttt{UOQ}, \texttt{UIQ}, \texttt{LOQ}, \texttt{LIQ}, \texttt{Central}) \\
\hline
14 & \texttt{s11\_deep\_margin} & ระยะห่างจากขอบตัดด้านลึก (\texttt{Deep Margin}, ซม.) --- \textbf{ฟิลด์วิกฤต (Critical)} \\
\hline
15 & \texttt{s14\_check} & การตรวจพบและระบุจำนวนต่อมน้ำเหลืองรักแร้ (\texttt{Axillary lymph nodes}) --- \textbf{ฟิลด์วิกฤต (Critical)} \\
\hline
\end{tabular}
\end{table}

\noindent ชุดข้อมูล \texttt{dataset\_ft} ได้รับการออกแบบให้มีข้อมูลครบทั้ง 15 ฟิลด์ครอบคลุมทั้ง 10 หมวดหมู่สภาวะทางคลินิก โดยตัวสกัดข้อมูลแบบกำหนดแน่นอนรุ่นที่ 3 (\texttt{src/nlp/extractor\_v3.py}) ให้ความถูกต้องเฉลี่ย $\text{Macro-}F_1 = 100.0\%$ ($\text{CER}_{\text{case}} = 0.00\%$) บนข้อความอ้างอิงเฉลยต้นฉบับ

\section*{ตารางแจงนับผลการสกัดข้อมูล (Confusion Matrix) รายฟิลด์บนชุดทดสอบรวม ($N = 300$ กรณี)}

\begin{table}[H]
\centering
\caption{จำนวนนับผลบวกจริง (TP), ผลบวกลวง (FP), ผลลบลวง (FN) และผลลบจริง (TN) ต่อฟิลด์บนชุดทดสอบรวม 300 กรณี (\text{test\_seen} 150 กรณี + \text{test\_unseen} 150 กรณี) เปรียบเทียบระหว่างระบบพื้นฐาน ($S_0$) และระบบบูรณาการ $S_3$ (กรณีสกัดค่าผิดไปจากเฉลยจะถูกนับเป็นทั้ง FP และ FN)}
\label{tab:confmat}
\footnotesize\setlength{\tabcolsep}{4pt}
\begin{tabular}{|l|c|c|c|c||c|c|c|c|}
\hline
 & \multicolumn{4}{c||}{\textbf{ระบบพื้นฐาน ($S_0$)}} & \multicolumn{4}{c|}{\textbf{ระบบบูรณาการ ($S_3$)}} \\
\hline
\textbf{ฟิลด์} & \textbf{TP} & \textbf{FP} & \textbf{FN} & \textbf{TN} & \textbf{TP} & \textbf{FP} & \textbf{FN} & \textbf{TN} \\
\hline
\texttt{s0\_surgical\_no} & 283 & 14 & 17 & 0 & 288 & 12 & 12 & 0 \\
\hline
\texttt{s1\_side} & 284 & 0 & 16 & 0 & 292 & 0 & 8 & 0 \\
\hline
\texttt{s2\_proc} & 290 & 0 & 10 & 0 & 296 & 0 & 4 & 0 \\
\hline
\texttt{s3\_dims} & 276 & 5 & 24 & 0 & 292 & 4 & 8 & 0 \\
\hline
\texttt{s4\_skin} & 146 & 0 & 4 & 150 & 147 & 0 & 3 & 150 \\
\hline
\texttt{s5\_dims} & 120 & 3 & 0 & 177 & 120 & 4 & 0 & 176 \\
\hline
\texttt{s6\_nipple} & 28 & 2 & 2 & 270 & 27 & 3 & 3 & 270 \\
\hline
\texttt{s7\_biopsy\_scar} & 28 & 0 & 2 & 270 & 30 & 0 & 0 & 270 \\
\hline
\texttt{s8\_cavity} & 30 & 0 & 0 & 270 & 30 & 0 & 0 & 270 \\
\hline
\texttt{s9\_residual\_mass} & 30 & 0 & 0 & 270 & 29 & 1 & 1 & 270 \\
\hline
\texttt{s10\_infiltrative} & 204 & 5 & 6 & 85 & 206 & 0 & 4 & 90 \\
\hline
\texttt{s10\_inf\_dims} & 197 & 6 & 13 & 90 & 202 & 4 & 8 & 90 \\
\hline
\texttt{s10\_5\_quadrant} & 148 & 2 & 2 & 150 & 147 & 2 & 3 & 150 \\
\hline
\texttt{s11\_deep\_margin} & 186 & 7 & 24 & 90 & 197 & 3 & 13 & 90 \\
\hline
\texttt{s14\_check} & 179 & 0 & 1 & 120 & 178 & 3 & 2 & 117 \\
\hline
\textbf{รวมทั้ง 15 ฟิลด์} & \textbf{2,429} & \textbf{44} & \textbf{121} & \textbf{1,942} & \textbf{2,481} & \textbf{36} & \textbf{69} & \textbf{1,943} \\
\hline
\end{tabular}
\end{table}

\noindent จากช่องข้อมูลทั้งหมด $300 \times 15 = 4{,}500$ ช่อง ระบบบูรณาการ $S_3$ ให้ความถูกต้องระดับช่องข้อมูล ($\text{TP} + \text{TN}$) สูงถึง $4{,}424 / 4{,}500$ ช่อง (คิดเป็นร้อยละ $98.31$) และลดจำนวนการตกหล่นหรือสกัดผิดพลาดของข้อมูลที่มีอยู่จริง (FN) ลงจาก $121$ ช่อง เหลือเพียง $69$ ช่อง (ลดลงร้อยละ $43.0$)

\section*{การกระจายตัวของข้อผิดพลาดตามหมวดหมู่การทดสอบ ($N = 300$ กรณี)}

\begin{table}[H]
\centering
\caption{จำนวนช่องข้อมูลที่สกัดผิดของระบบบูรณาการ $S_3$ จำแนกตามฟิลด์และหมวดหมู่การทดสอบทั้ง 10 หมวดหมู่ (หมวดหมู่ละ 30 กรณี: \text{test\_seen} 15 กรณี + \text{test\_unseen} 15 กรณี)}
\label{tab:caterr}
\footnotesize\setlength{\tabcolsep}{4pt}
\begin{tabular}{|l|c|c|c|c|c|c|c|c|c|c|c|}
\hline
\textbf{ฟิลด์} & \textbf{1} & \textbf{2} & \textbf{3} & \textbf{4} & \textbf{5} & \textbf{6} & \textbf{7} & \textbf{8} & \textbf{9} & \textbf{10} & \textbf{รวม} \\
\hline
\texttt{s0\_surgical\_no} & 0 & 0 & 0 & 1 & 0 & 0 & 2 & 9 & 0 & 0 & 12 \\
\hline
\texttt{s1\_side} & 0 & 0 & 0 & 0 & 2 & 0 & 0 & 6 & 0 & 0 & 8 \\
\hline
\texttt{s2\_proc} & 0 & 0 & 0 & 0 & 2 & 0 & 0 & 2 & 0 & 0 & 4 \\
\hline
\texttt{s3\_dims} & 0 & 0 & 0 & 0 & 2 & 0 & 0 & 6 & 0 & 0 & 8 \\
\hline
\texttt{s4\_skin} & 0 & 0 & 3 & 0 & 0 & 0 & 0 & 0 & 0 & 0 & 3 \\
\hline
\texttt{s5\_dims} & 0 & 0 & 4 & 0 & 0 & 0 & 0 & 0 & 0 & 0 & 4 \\
\hline
\texttt{s6\_nipple} & 0 & 0 & 0 & 3 & 0 & 0 & 0 & 0 & 0 & 0 & 3 \\
\hline
\texttt{s7\_biopsy\_scar} & 0 & 0 & 0 & 0 & 0 & 0 & 0 & 0 & 0 & 0 & 0 \\
\hline
\texttt{s8\_cavity} & 0 & 0 & 0 & 0 & 0 & 0 & 0 & 0 & 0 & 0 & 0 \\
\hline
\texttt{s9\_residual\_mass} & 0 & 0 & 0 & 0 & 0 & 0 & 0 & 0 & 0 & 1 & 1 \\
\hline
\texttt{s10\_infiltrative} & 0 & 0 & 0 & 0 & 0 & 0 & 0 & 4 & 0 & 0 & 4 \\
\hline
\texttt{s10\_inf\_dims} & 0 & 0 & 1 & 0 & 0 & 0 & 1 & 6 & 0 & 0 & 8 \\
\hline
\texttt{s10\_5\_quadrant} & 1 & 0 & 0 & 1 & 0 & 0 & 0 & 0 & 1 & 0 & 3 \\
\hline
\texttt{s11\_deep\_margin} & 0 & 0 & 0 & 0 & 0 & 0 & 0 & 12 & 0 & 1 & 13 \\
\hline
\texttt{s14\_check} & 0 & 0 & 0 & 3 & 2 & 0 & 0 & 0 & 0 & 0 & 5 \\
\hline
\textbf{รวมช่องที่ผิด} & \textbf{1} & \textbf{0} & \textbf{8} & \textbf{8} & \textbf{8} & \textbf{0} & \textbf{3} & \textbf{45} & \textbf{1} & \textbf{2} & \textbf{76} \\
\hline
\end{tabular}
\end{table}

\noindent จากตารางจะเห็นได้ว่าในหมวดหมู่ที่ 2 (การบรรยายข้อมูลสลับลำดับ: Out-of-Order Reporting) และหมวดหมู่ที่ 6 (สิ่งส่งตรวจเนื้อเยื่อปกติที่ไม่ใช่มะเร็ง: Benign Control) ระบบบูรณาการ $S_3$ มีความผิดพลาดเป็นศูนย์ ($0$ ช่อง) ส่วนความผิดพลาดส่วนใหญ่กระจุกตัวอยู่ในหมวดหมู่ที่ 8 (สภาวะเสียงรบกวนพัดลมตู้ดูดควันระดับรุนแรง $\mathrm{SNR} = -10\text{ dB}$) ซึ่งระบบแจ้งเตือนความปลอดภัยหลายสัญญาณจะยกธงแจ้งเตือนความเสี่ยงในสัดส่วนสูงที่สุดเพื่อส่งต่อให้พยาธิแพทย์ตรวจทานซ้ำตามระเบียบวิธีคัดกรองความปลอดภัย

\section*{การทดสอบระดับหน่วย (Unit Testing) ของตัวปรับข้อความและตัวสกัดข้อมูล}

\begin{table}[H]
\centering
\caption{ผลการทดสอบระดับหน่วยของตัวปรับข้อความสำหรับคัดกรองการพูดแก้ไขตัวเองกลางประโยค (16 กรณีทดสอบ) และการทดสอบบูรณาการระดับระบบ (4 ชุดการทดสอบหลัก)}
\label{tab:unit}
\footnotesize\setlength{\tabcolsep}{3.5pt}
\begin{tabular}{|p{1.7cm}|p{5.0cm}|p{5.0cm}|p{1.4cm}|}
\hline
\textbf{รหัส} & \textbf{ข้อมูลนำเข้า / ชุดการทดสอบ} & \textbf{ผลลัพธ์ที่ได้} & \textbf{สถานะ} \\
\hline
TC-01 & \texttt{infiltrative mass measuring 2.0 x 3.0 x 1.0 cm sorry 2.5 x 3.5 x 1.5 cm} & \texttt{infiltrative mass measuring 2.5 x 3.5 x 1.5 cm} & ผ่าน \\
\hline
TC-02..06 & การแยกแยะคำว่า \texttt{weight / grams / kg} ออกจากคำกริยาแจ้งแก้ไข \texttt{wait} & คงค่าขนาดชิ้นเนื้อและน้ำหนักไว้ครบถ้วนโดยไม่ถูกลบ & ผ่าน \\
\hline
TC-07..09 & การใช้คำว่า \texttt{actually / no} ในบริบทเปลี่ยนหัวข้อฟิลด์หรือระบุจำนวนต่อมน้ำเหลือง & ไม่ลบค่าของฟิลด์ก่อนหน้าโดยผิดพลาด & ผ่าน \\
\hline
TC-10..11 & \texttt{mass measuring 2.0 x 3.0 x 1.0 cm weight (wait) 2.5 x 3.5 x 1.5 cm} & \texttt{mass measuring 2.5 x 3.5 x 1.5 cm} & ผ่าน \\
\hline
TC-12..15 & ประโยคปฏิเสธทางคลินิก (\texttt{no residual mass / no discrete mass}) & คงประโยคปฏิเสธไว้ครบถ้วนและตั้งค่า \texttt{No mass} ถูกต้อง & ผ่าน \\
\hline
TC-16 & \texttt{mass measuring 2 x 2 cm แก้เป็น 3 x 3 cm} & \texttt{mass measuring 3 x 3 cm} & ผ่าน \\
\hline
E2E-01 & ทดสอบการกำหนดค่า \texttt{config.py}, แบบจำลอง LoRA INT8 และข้อความนำ $P_4$ & โหลดแบบจำลอง \texttt{models/pathowhisper\_lora\_int8} สำเร็จ & ผ่าน \\
\hline
E2E-02 & ทดสอบ \texttt{extractor\_v3.py} บนชุดข้อความตัวแทนครบทั้ง 15 ฟิลด์ & สกัดข้อมูลครบ 15/15 ฟิลด์ด้วยความถูกต้องร้อยละ $100$ & ผ่าน \\
\hline
E2E-03 & ทดสอบเอนจิน \texttt{FasterWhisperEngine} ร่วมกับกลไก \texttt{D8\_snr\_gated} & คัดแยกเสียงสะอาด ($\widehat{\mathrm{SNR}} = 49.4\text{ dB}$) และเสียงตู้ดูดควัน ($\widehat{\mathrm{SNR}} = 1.9\text{ dB}$) ถูกต้อง & ผ่าน \\
\hline
E2E-04 & ทดสอบจุดเชื่อมต่อ \texttt{POST /api/transcribe} ของเว็บแอปพลิเคชัน Flask แบบครบวงจร & ถอดความ สกัดข้อมูล 15 ฟิลด์ และสร้างธงเตือนสำเร็จในเวลา $3.40\text{ s}$ & ผ่าน \\
\hline
\end{tabular}
\end{table}

\appendiceschapter{ข้อมูลและสคริปต์ที่ใช้ในการทดลองและคำนวณผล}
\section*{คลังข้อมูล สคริปต์การทดลอง และไฟล์ผลลัพธ์}

ซอร์สโค้ดของโครงงาน สคริปต์การทดลอง ชุดข้อมูลควบคุม (Manifests และ Ground Truths) ตลอดจนไฟล์ผลลัพธ์การคำนวณสถิติและรูปภาพประกอบวิทยานิพนธ์ ได้รับการเผยแพร่แบบโอเพนซอร์สบนคลังรหัสต้นฉบับ GitHub เพื่อความโปร่งใสและการทำซ้ำผลการวิจัย (Reproducibility):
\begin{center}
\textbf{GitHub Repository:} \url{https://github.com/roojask/app_pathology_7-7-2569}
\end{center}

โครงสร้างของคลังข้อมูลและสคริปต์ได้รับการออกแบบตามหลักการแยกส่วนอย่างเคร่งครัด (Modular Architecture) เพื่อให้สามารถตรวจสอบย้อนกลับ (Auditability) และทำซ้ำการทดลองได้ในทุกระดับ โดยแบ่งออกเป็น 5 หมวดหมู่หลักตามตารางที่ \ref{tab:appfiles}

{\footnotesize\setlength{\tabcolsep}{3.5pt}
\begin{longtable}{|>{\raggedright\arraybackslash}p{5.6cm}|p{9.4cm}|}
\caption{โครงสร้างไฟล์ข้อมูล สคริปต์การทดลองแบบแยกส่วน และไฟล์ผลลัพธ์ในสถาปัตยกรรมการทดลอง}\label{tab:appfiles}\\
\hline
\textbf{ไฟล์ / โฟลเดอร์} & \textbf{หน้าที่และเนื้อหา} \\
\hline
\endfirsthead
\caption[]{โครงสร้างไฟล์ข้อมูล สคริปต์การทดลองแบบแยกส่วน และไฟล์ผลลัพธ์ในสถาปัตยกรรมการทดลอง (ต่อ)}\\
\hline
\textbf{ไฟล์ / โฟลเดอร์} & \textbf{หน้าที่และเนื้อหา} \\
\hline
\endhead
\multicolumn{2}{|l|}{\textbf{1. ชุดข้อมูลและข้อมูลเฉลยมาตรฐาน (Synthetic Dataset and Ground Truth)}} \\
\hline
\texttt{scripts/\allowbreak{}make\_\allowbreak{}finetune\_\allowbreak{}data.py} & สคริปต์สร้างชุดข้อมูลสังเคราะห์ 1,000 กรณี พร้อมแยกคลังสำนวนประโยคระหว่าง \texttt{train} (600), \texttt{val} (100), \texttt{test\_\allowbreak{}seen} (150) และ \texttt{test\_\allowbreak{}unseen} (150) ด้วยหลักการ Zero-Phrase-Leakage และทดสอบทวนสอบตัวสกัด 15 ฟิลด์ได้ค่า Macro-$F_1 = 100\%$ \\
\hline
\texttt{data/\allowbreak{}dataset\_\allowbreak{}ft/\allowbreak{}ground\_\allowbreak{}truth\_\allowbreak{}ft.json}, \texttt{manifest.csv} & ไฟล์ข้อมูลเฉลยครบทั้ง 15 ฟิลด์ CAP และรายการดัชนีไฟล์เสียงความละเอียดสูง 16 kHz ทั้ง 1,000 กรณีศึกษา \\
\hline
\texttt{benchmarks/\allowbreak{}clean\_\allowbreak{}slate/\allowbreak{}results/\allowbreak{}calibrated\_\allowbreak{}tau.json} & ค่าเกณฑ์ขีดแบ่งอัตราส่วนสัญญาณต่อเสียงรบกวน ($\tau^* = 13.55\text{ dB}$) ที่สอบเทียบจากดัชนียูเดนบนชุดตรวจสอบ \\
\hline
\multicolumn{2}{|l|}{\textbf{2. โมดูลระบบหลักและตัวสกัดข้อมูล (Core Pipeline and Parser Modules)}} \\
\hline
\texttt{src/\allowbreak{}nlp/\allowbreak{}extractor\_\allowbreak{}v3.py} & ตัวสกัดข้อมูลแบบกำหนดแน่นอนรุ่นที่ 3 (Deterministic Rule-Based Parser) ครอบคลุม 15 ฟิลด์ CAP พร้อมระบบสร้างธงเตือนความปลอดภัยหลายสัญญาณ (Multiple Safety Risk Flags) \\
\hline
\texttt{src/\allowbreak{}core/\allowbreak{}engine\_\allowbreak{}faster\_\allowbreak{}whisper.py} & โมดูลเอนจินการถอดรหัสเสียงออฟไลน์ เชื่อมต่อโมเดล CTranslate2 INT8, กลไกคัดกรองสัญญาณ \texttt{D8\_\allowbreak{}snr\_\allowbreak{}gated} และพารามิเตอร์การถอดความ \\
\hline
\texttt{config.py} & ไฟล์กำหนดค่าพารามิเตอร์ระบบส่วนกลาง (พาธแบบจำลอง, ค่าขีดแบ่ง $\tau^*$, ข้อความนำ $P_4$, โฮสต์และพอร์ตการทำงาน) \\
\hline
\multicolumn{2}{|l|}{\textbf{3. ชุดสคริปต์การทดลองและการคำนวณตัวชี้วัด (Clean-Slate Benchmark Suite)}} \\
\hline
\texttt{benchmarks/\allowbreak{}clean\_\allowbreak{}slate/\allowbreak{}acoustic\_\allowbreak{}filters.py} & โมดูลคำนวณค่าประมาณ $\widehat{\mathrm{SNR}}$ แบบไม่ต้องใช้สัญญาณอ้างอิง และกลไกคัดกรองสัญญาณเสียงแบบปรับตัว \texttt{D8\_\allowbreak{}snr\_\allowbreak{}gated} \\
\hline
\texttt{benchmarks/\allowbreak{}clean\_\allowbreak{}slate/\allowbreak{}metrics\_\allowbreak{}clean.py} & โมดูลคำนวณตัวชี้วัดทางการแพทย์และ ASR (Normalized WER, C-TER, N-DER, 15-Field Macro-$F_1$, $\mathrm{CER}_{\mathrm{case}}$ และ Clinical Risk Score) \\
\hline
\texttt{benchmarks/\allowbreak{}clean\_\allowbreak{}slate/\allowbreak{}run\_\allowbreak{}clean\_\allowbreak{}slate\_\allowbreak{}experiments.py} & สคริปต์หลักสำหรับรันการทดลองที่ 1--3 (การสอบเทียบ $\tau^*$, ตารางกริดค้นหาพารามิเตอร์ \texttt{afftdn}, การแยกส่วนข้อความนำ $P_0 \dots P_4$ และการทดสอบโมเดลฐาน) \\
\hline
\texttt{benchmarks/\allowbreak{}clean\_\allowbreak{}slate/\allowbreak{}expand\_\allowbreak{}arch\_\allowbreak{}study.py} & สคริปต์ประเมินเปรียบเทียบข้ามสถาปัตยกรรม (Kaldi FST Vosk, Wav2Vec 2.0 Base/Large CTC, HuBERT, WavLM และ Whisper Base/Small) \\
\hline
\texttt{benchmarks/\allowbreak{}clean\_\allowbreak{}slate/\allowbreak{}train\_\allowbreak{}whisper\_\allowbreak{}lora.py} & สคริปต์ฝึกฝน LoRA ($r=8, \alpha=32$) บนชุดข้อมูล \texttt{train} 600 กรณี, การแปลงแบบจำลองเป็น CTranslate2 INT8 และประเมินผลการลดความผิดพลาดตามลำดับขั้น ($S_0 \dots S_3$) \\
\hline
\texttt{benchmarks/\allowbreak{}clean\_\allowbreak{}slate/\allowbreak{}fix\_\allowbreak{}longform\_\allowbreak{}lora\_\allowbreak{}eval.py} & สคริปต์ประมวลผลการถอดรหัสตามความยาวเสียง (Duration-Aware Long-Form Decoding) สำหรับไฟล์เสียงที่ยาวเกิน 27.5 วินาที \\
\hline
\multicolumn{2}{|l|}{\textbf{4. การวิเคราะห์เชิงสถิติ ตารางสรุป และรูปภาพประกอบ (Statistical Analysis and Results)}} \\
\hline
\texttt{benchmarks/\allowbreak{}clean\_\allowbreak{}slate/\allowbreak{}generate\_\allowbreak{}final\_\allowbreak{}thesis\_\allowbreak{}tables.py} & สคริปต์รวบรวมผลลัพธ์การทดลองทั้งหมด คำนวณช่วงความเชื่อมั่น 95\% ทดสอบนัยสำคัญทางสถิติ (McNemar's test, Wilcoxon signed-rank test พร้อมการปรับแก้แบบ Holm-Bonferroni) และสร้างตารางสรุปผลวิทยานิพนธ์ (ตารางที่ 6--13) \\
\hline
\texttt{benchmarks/\allowbreak{}clean\_\allowbreak{}slate/\allowbreak{}generate\_\allowbreak{}thesis\_\allowbreak{}figures.py} & สคริปต์พล็อตและเรนเดอร์ภาพกราฟประกอบวิทยานิพนธ์ทั้ง 6 ภาพ (ภาพที่ 1--6) ที่ความละเอียดสูง 300 DPI ด้วย Matplotlib และ Seaborn \\
\hline
\texttt{benchmarks/\allowbreak{}clean\_\allowbreak{}slate/\allowbreak{}results/\allowbreak{}final\_\allowbreak{}thesis\_\allowbreak{}master\_\allowbreak{}summary.json} & ไฟล์ฐานข้อมูลผลการทดลองฉบับสมบูรณ์ (Master Summary JSON) ที่รวบรวมค่าตัวชี้วัดทุกระดับของระบบ $S_0 \dots S_3$ ไว้อย่างเป็นเอกภาพ \\
\hline
\texttt{benchmarks/\allowbreak{}clean\_\allowbreak{}slate/\allowbreak{}results/\allowbreak{}*.csv} & ไฟล์บันทึกผลลัพธ์รายกรณี ($N=150$) และตารางสรุปของการทดลองทั้งหมด (\texttt{snr\_\allowbreak{}calibration\_\allowbreak{}*}, \texttt{exp1\_\allowbreak{}*}, \texttt{exp2\_\allowbreak{}*}, \texttt{exp3\_\allowbreak{}*}, \texttt{exp4\_\allowbreak{}*}, \texttt{final\_\allowbreak{}table\_\allowbreak{}*}) \\
\hline
\multicolumn{2}{|l|}{\textbf{5. แอปพลิเคชันและตัวติดตั้งสำหรับใช้งานจริง (Application and Desktop Deployment)}} \\
\hline
\texttt{app.py} & เซิร์ฟเวอร์ Flask ให้บริการจุดเชื่อมต่อ API ถอดความเสียง สกัดข้อมูล และหน้าอินเทอร์เฟซผู้ใช้ \\
\hline
\texttt{gui\_\allowbreak{}app.py} & แอปพลิเคชัน Desktop พร้อมถาดระบบ (System Tray) สำหรับเปิดปิดบริการเซิร์ฟเวอร์ภายในเครื่องแบบ Standalone \\
\hline
\texttt{tests/\allowbreak{}test\_\allowbreak{}app\_\allowbreak{}clean\_\allowbreak{}slate\_\allowbreak{}e2e.py} & สคริปต์ทดสอบการทำงานร่วมกันแบบครบวงจร (End-to-End System Tests) ของระบบบูรณาการ $S_3$ \\
\hline
\texttt{build\_\allowbreak{}installer.ps1}, \texttt{installer.iss} & สคริปต์ PowerShell และ Inno Setup สำหรับสร้างแพ็กเกจตัวติดตั้งแบบ Standalone ออฟไลน์บนระบบปฏิบัติการ Windows \\
\hline
\end{longtable}}

\section*{คำอธิบายโครงสร้างไฟล์ผลลัพธ์การทดลอง (Experiment Output Data Structure)}

ไฟล์ผลลัพธ์ในไดเรกทอรี \texttt{benchmarks/clean\_slate/results/} จัดเก็บข้อมูลการทดลองอย่างเป็นระบบใน 3 ระดับ:
\begin{enumerate}
    \item \textbf{ไฟล์ผลลัพธ์ระดับรายกรณี (Case-Level CSVs):} เช่น \texttt{exp4\_lora\_evaluation\_test\_seen.csv} และ \texttt{exp4\_lora\_evaluation\_test\_unseen.csv} บันทึกข้อความถอดความจริง ผลการทำนายของตัวสกัดทั้ง 15 ฟิลด์ ค่าความผิดพลาดของคำศัพท์เฉพาะ (C-TER) และตัวเลข (N-DER) เทียบกับเฉลยสำหรับทุกกรณีศึกษา ($N = 150$) เพื่อใช้ในการตรวจสอบย้อนกลับและวิเคราะห์ข้อผิดพลาด
    \item \textbf{ไฟล์ผลลัพธ์ระดับตารางสรุป (Summary CSVs):} เช่น \texttt{final\_table\_system\_progression\_test\_seen.csv}, \texttt{final\_table\_prompt\_ablation\_test\_seen.csv} และ \texttt{final\_table\_15fields\_test\_seen.csv} ซึ่งเก็บค่าเฉลี่ยและส่วนเบี่ยงเบนมาตรฐานที่ตรงกับตารางในบทที่ 4 ทุกประการ
    \item \textbf{ไฟล์ฐานข้อมูลสถิติรวมศูนย์ (\texttt{final\_thesis\_master\_summary.json}):} ทำหน้าที่เป็น Single Source of Truth บันทึกสถิติสรุป ค่าสถิติอนุมาน (ค่าสถิติทดสอบ $W$ และ $\chi^2$, ค่า $p$ ก่อนและหลังปรับแก้ Holm-Bonferroni) และช่วงความเชื่อมั่นร้อยละ 95 เพื่อให้สคริปต์สร้างตารางและสร้างรูปภาพดึงข้อมูลจากแหล่งเดียวกันอย่างสอดคล้องสมบูรณ์
\end{enumerate}


\appendiceschapter{หลักฐานการทดสอบการทำงานแบบออฟไลน์}

การทดสอบระบบในสภาวะตัดการเชื่อมต่อเครือข่ายภายนอก (Isolated Environment Testing) ดำเนินการบนเครื่องคอมพิวเตอร์ส่วนบุคคลที่ใช้ในการทดลอง เพื่อยืนยันว่าระบบ PathoWhisper สามารถประมวลผลสัญญาณเสียง สกัดข้อมูลลงแบบฟอร์ม และสร้างรายงานสรุปในรูปแบบเอกสาร PDF ได้อย่างสมบูรณ์โดยไม่พึ่งพาการเชื่อมต่ออินเทอร์เน็ตหรือบริการคลาวด์ภายนอก ซึ่งสอดคล้องกับมาตรการคุ้มครองข้อมูลส่วนบุคคลทางการแพทย์ ขั้นตอนและผลการทดสอบเชิงฟังก์ชันสรุปดังตารางที่ \ref{tab:offline}

\begin{table}[H]
\centering
\caption{สรุปขั้นตอนและผลการทดสอบการทำงานของระบบในสภาวะออฟไลน์}
\label{tab:offline}
\small
\begin{tabular}{|c|p{4.2cm}|p{7.8cm}|}
\hline
\textbf{ลำดับ} & \textbf{ขั้นตอนการทดสอบ} & \textbf{ผลการทดสอบและหลักฐานเชิงประจักษ์} \\
\hline
1 & การตรวจสอบสถานะเครือข่าย (ภาพที่ \ref{fig:off1}) & คำสั่ง \texttt{ping 8.8.8.8} ล้มเหลวทั้ง 4 ครั้ง (ส่ง 4 รับ 0 สูญหายร้อยละ 100) และสัญลักษณ์เครือข่ายบนแถบงานแสดงสถานะไม่มีการเชื่อมต่ออินเทอร์เน็ต \\
\hline
2 & การประมวลผลเสียงและสกัดข้อมูล (ภาพที่ \ref{fig:off2}) & ระบบถอดความสัญญาณเสียงความยาว 23 วินาที และสกัดข้อมูลลงแบบฟอร์มได้ครบถ้วน เช่น เลขที่สิ่งส่งตรวจ \texttt{S-24-1027}, ประเภทหัตถการ MRM ข้างขวา, ขนาดสิ่งส่งตรวจ และตำแหน่งก้อน \\
\hline
3 & การตรวจสอบเงื่อนไขความถูกต้อง (ภาพที่ \ref{fig:off3}) & เมื่อทดสอบด้วยกรณีที่ไม่มีการระบุเลขที่สิ่งส่งตรวจ ระบบแสดงข้อความแจ้งเตือนและเว้นว่างฟิลด์ดังกล่าวตามเงื่อนไขความปลอดภัยของข้อมูล \\
\hline
4 & การสร้างรายงานสรุปผล PDF (ภาพที่ \ref{fig:off4}) & ไฟล์เอกสารรายงาน \texttt{pathology\_breast\_gross\_S-24-1027 (1).pdf} ของกรณีศึกษาเดียวกับภาพที่ \ref{fig:off2} เปิดแสดงผลขณะที่ระบบยังคงตัดการเชื่อมต่อเครือข่าย โดยข้อมูลตรงกับเสียงบรรยายทุกประการ \\
\hline
\end{tabular}
\end{table}

\begin{figure}[H]
\centering
\includegraphics[width=0.97\textwidth]{images/evidence/offline_1_ping_1553.png}
\caption{การตรวจสอบสถานะการตัดการเชื่อมต่อเครือข่ายภายนอก (Network Isolation Verification)}
\label{fig:off1}
\end{figure}

\begin{figure}[H]
\centering
\includegraphics[width=0.97\textwidth]{images/evidence/offline_4_pdf_1742.png}
\caption{การประมวลผลถอดความเสียงและสกัดข้อมูลลงแบบฟอร์มในสภาวะออฟไลน์}
\label{fig:off2}
\end{figure}

\begin{figure}[H]
\centering
\includegraphics[width=0.97\textwidth]{images/evidence/offline_3_missing_field_flag_1600.png}
\caption{การแสดงข้อความแจ้งเตือนของระบบเมื่อไม่พบข้อมูลเลขที่สิ่งส่งตรวจทางศัลยกรรม}
\label{fig:off3}
\end{figure}

\begin{figure}[H]
\centering
\includegraphics[width=0.97\textwidth]{images/evidence/offline_2_dictation_1742.png}
\caption{เอกสารรายงาน PDF ของกรณีศึกษา S-24-1027 ที่เปิดแสดงผลขณะตัดการเชื่อมต่อเครือข่าย}
\label{fig:off4}
\end{figure}

\noindent \textbf{ข้อสังเกตเชิงเทคนิคจากการพัฒนาและปรับปรุงระบบ:}
\begin{itemize}
    \item \textbf{การขยายความครอบคลุมครบทั้ง 15 ฟิลด์:} ในตัวสกัดข้อมูลรุ่นที่ 3 (\texttt{extractor\_v3.py}) ได้เพิ่มกฎการสกัดข้อมูลสำหรับฟิลด์สภาพหัวนม (\texttt{s6\_nipple}), รอยแผลเป็นจากการเจาะตรวจ (\texttt{s7\_biopsy\_scar}), โพรงผ่าตัดเดิม (\texttt{s8\_cavity}) และก้อนเนื้องอกหลงเหลือ (\texttt{s9\_residual\_mass}) จนครบสมบูรณ์ทั้ง 15 ฟิลด์ตามมาตรฐาน CAP
    \item \textbf{การแยกแยะขนาดชิ้นผิวหนังกับระยะห่างขอบตัดผ่าตัด:} ตัวสกัดใช้บริบทคำว่า \texttt{margin} กำกับระยะห่างขอบตัดเสมอ จึงไม่เกิดการสับสนระหว่างขนาดวงรีผิวหนัง (\texttt{s5\_dims}) กับระยะห่างจากขอบตัดด้านลึก (\texttt{s11\_deep\_margin})
    \item \textbf{การจัดสรรชุดตัวเลขมิติตามบริบท:} แม้ผู้พูดจะสลับลำดับการบรรยายโดยระบุขนาดก้อนเนื้องอกก่อนขนาดสิ่งส่งตรวจหลัก ตัวสกัดข้อมูลสามารถวิเคราะห์คำบ่งชี้บริบทและจัดสรรค่ามิติลงช่อง \texttt{s10\_inf\_dims} และ \texttt{s3\_dims} ได้อย่างถูกต้อง
\end{itemize}

\appendiceschapter{คู่มือการติดตั้งแอปพลิเคชัน}
\label{app:setup install}


คู่มือฉบับนี้นำเสนอขั้นตอนการดาวน์โหลดและการติดตั้งแอปพลิเคชัน \textbf{Pathology Voice Assistant} ลงบนเครื่องคอมพิวเตอร์ส่วนบุคคล (PC) สำหรับสถานีตรวจชิ้นเนื้อทางศัลยพยาธิวิทยา โดยระบบได้รับการออกแบบเป็นตัวติดตั้งแบบอัตโนมัติ (Standard Setup Installer) เพื่ออำนวยความสะดวกให้ผู้ใช้งานและผู้ดูแลระบบสามารถติดตั้งและพร้อมใช้งานออฟไลน์ได้โดยง่าย

\section*{1. ข้อกำหนดขั้นต่ำของระบบ (System Requirements)}
ก่อนเริ่มดำเนินกระบวนการติดตั้ง ผู้ดูแลระบบหรือผู้ใช้งานควรตรวจสอบคุณสมบัติของเครื่องคอมพิวเตอร์ให้ตรงตามข้อกำหนดขั้นต่ำของระบบ ดังนี้:
\begin{itemize}
    \item \textbf{ระบบปฏิบัติการ (Operating System):} Microsoft Windows 10 หรือ Windows 11 (64-bit)
    \item \textbf{หน่วยประมวลผลกลาง (CPU):} Intel Core i5 / AMD Ryzen 5 ขึ้นไป (แนะนำหน่วยประมวลผลขนาด 8 เธรดขึ้นไป)
    \item \textbf{หน่วยความจำหลัก (RAM):} 8 GB ขึ้นไป (แนะนำ 16 GB เพื่อประสิทธิภาพการถอดความสูงสุด)
    \item \textbf{พื้นที่จัดเก็บข้อมูล (Hard Disk Space):} พื้นที่ว่างอย่างน้อย 5 GB (สำหรับตัวโปรแกรม ค่าน้ำหนักแบบจำลองออฟไลน์ และฐานข้อมูลภายในเครื่อง)
    \item \textbf{อุปกรณ์เชื่อมต่อภายนอก (Peripherals):} ไมโครโฟนรับสัญญาณเสียง และแป้นเหยียบสั่งงานผ่านพอร์ต USB (USB Foot Pedal)
\end{itemize}

\section*{2. ขั้นตอนการดาวน์โหลดและติดตั้งแอปพลิเคชัน}

\section*{1. ข้อกำหนดขั้นต่ำของระบบ (System Requirements)}
ก่อนเริ่มดำเนินกระบวนการติดตั้ง ผู้ดูแลระบบหรือผู้ใช้งานควรตรวจสอบคุณสมบัติของเครื่องคอมพิวเตอร์ให้ตรงตามข้อกำหนดขั้นต่ำของระบบ ดังนี้:
\begin{itemize}
    \item \textbf{ระบบปฏิบัติการ (Operating System):} Microsoft Windows 10 หรือ Windows 11 (64-bit)
    \item \textbf{หน่วยประมวลผลกลาง (CPU):} Intel Core i5 / AMD Ryzen 5 ขึ้นไป (แนะนำหน่วยประมวลผลขนาด 8 เธรดขึ้นไป)
    \item \textbf{หน่วยความจำหลัก (RAM):} 8 GB ขึ้นไป (แนะนำ 16 GB เพื่อประสิทธิภาพการถอดความสูงสุด)
    \item \textbf{พื้นที่จัดเก็บข้อมูล (Hard Disk Space):} พื้นที่ว่างอย่างน้อย 5 GB (สำหรับตัวโปรแกรม ค่าน้ำหนักแบบจำลองออฟไลน์ และฐานข้อมูลภายในเครื่อง)
    \item \textbf{อุปกรณ์เชื่อมต่อภายนอก (Peripherals):} ไมโครโฟนรับสัญญาณเสียง และแป้นเหยียบสั่งงานผ่านพอร์ต USB (USB Foot Pedal)
\end{itemize}

\section*{2. ขั้นตอนการดาวน์โหลดและติดตั้งแอปพลิเคชัน}

\subsection*{2.1 ขั้นตอนที่ 1: การดาวน์โหลดไฟล์ติดตั้งจาก Google Drive}
\begin{enumerate}
    \item[2.1.1] เข้าสู่ระบบผ่านลิงก์เชื่อมโยง Google Drive ที่ได้รับมอบหมายเพื่อดาวน์โหลดชุดไฟล์ติดตั้งโปรแกรม
    \item[2.1.2] คลิกปุ่ม \textbf{"ดาวน์โหลด" (Download)} ไฟล์ติดตั้งแอปพลิเคชัน ชื่อ \texttt{PathologyApp\_Setup.exe} ลงในเครื่องคอมพิวเตอร์
\end{enumerate}

\begin{figure}[htbp]
    \centering
    \includegraphics[width=0.8\textwidth]{คู่มือการติดตั้ง/Screenshot 2026-10-22 161359.png}
    \caption{แสดงหน้าต่างการดาวน์โหลดไฟล์ติดตั้ง \texttt{PathologyApp\_Setup.exe} จากระบบ Google Drive}
    \label{fig:fig1}
\end{figure}

\subsection*{2.2 ขั้นตอนที่ 2: การเปิดรันไฟล์ติดตั้งโปรแกรม (Installer Execution)}
\begin{enumerate}
    \item[2.2.1] ดับเบิลคลิกที่ไฟล์ \texttt{PathologyApp\_Setup.exe} เพื่อเข้าสู่โปรแกรมช่วยติดตั้ง (Setup Wizard)
    \item[2.2.2] หากระบบปฏิบัติการแสดงหน้าต่างแจ้งเตือนสิทธิ์การใช้งาน (User Account Control: UAC) ให้คลิกตอบรับ \textbf{"Yes"} หรือ \textbf{"Run anyway"} เพื่อยินยอมให้โปรแกรมดำเนินการติดตั้ง
\end{enumerate}

\begin{figure}[htbp]
    \centering
    \includegraphics[width=0.8\textwidth]{คู่มือการติดตั้ง/Screenshot 2026-10-22 162642.png}
    \caption{แสดงหน้าต่างการแจ้งเตือนสิทธิ์การใช้งาน (UAC) และการกดอนุญาตเพื่อรันโปรแกรมติดตั้ง}
    \label{fig:fig2}
\end{figure}

\subsection*{2.3 ขั้นตอนที่ 3: ตำแหน่งโฟลเดอร์ปลายทางสำหรับการติดตั้ง}
\begin{enumerate}
    \item[2.3.1] หน้าต่างตั้งค่าจะเลือกตำแหน่งโฟลเดอร์ปลายทางมาตรฐานที่ \texttt{C:\textbackslash Program Files\textbackslash PathologyApp} โดยอัตโนมัติ
    \item[2.3.2] จากนั้นคลิกปุ่ม \textbf{"Next >"}
\end{enumerate}

\begin{figure}[htbp]
    \centering
    \includegraphics[width=0.8\textwidth]{คู่มือการติดตั้ง/Screenshot 2026-10-22 162742.png}
    \caption{หน้าต่างจะเลือกตำแหน่งโฟลเดอร์สำหรับจัดเก็บไฟล์โปรแกรม PathologyApp โดยอัตโนมัติ}
    \label{fig:fig3}
\end{figure}

\subsection*{2.4 ขั้นตอนที่ 4: การดำเนินการติดตั้งระบบและค่าน้ำหนักแบบจำลอง}
\begin{enumerate}
    \item[2.4.1] ตรวจสอบความถูกต้องของรายการตั้งค่า แล้วคลิกปุ่ม \textbf{"Install"}
    \item[2.4.2] โปรแกรมช่วยติดตั้งจะทำการคัดลอกไฟล์ซอร์สโค้ด ค่าน้ำหนักแบบจำลองออฟไลน์ และไฟล์ส่วนประกอบระบบลงบนเครื่องคอมพิวเตอร์
\end{enumerate}

\begin{figure}[htbp]
    \centering
    \includegraphics[width=0.8\textwidth]{คู่มือการติดตั้ง/Screenshot 2026-10-22 162925.png}
    \caption{แสดงหน้าต่างแถบแสดงความก้าวหน้า (Progress Bar) ระหว่างการคัดลอกไฟล์ลงในระบบ}
    \label{fig:fig4}
\end{figure}

\subsection*{2.5 ขั้นตอนที่ 5: การเสร็จสิ้นการติดตั้งและการเปิดใช้งานครั้งแรก}
\begin{enumerate}
    \item[2.5.1] เมื่อคัดลอกไฟล์เสร็จสมบูรณ์ ทำเครื่องหมายหน้าช่อง \textbf{"Launch PathologyApp"} แล้วคลิกปุ่ม \textbf{"Finish"}
    \item[2.5.2] ระบบจะสร้างไอคอนทางลัด (Shortcut Icon) ชื่อ \textbf{PathologyApp} ไว้บนหน้าจอเดสก์ท็อปเพื่อความสะดวกในการเปิดใช้งาน
\end{enumerate}

\begin{figure}[htbp]
    \centering
    \includegraphics[width=0.8\textwidth]{คู่มือการติดตั้ง/Screenshot 2026-10-22 163128.png}
    \caption{แสดงหน้าต่างยืนยันการติดตั้งโปรแกรมเสร็จสมบูรณ์ และไอคอนทางลัดบนหน้าจอเดสก์ท็อป}
    \label{fig:fig5}
\end{figure}

\section*{3. การตรวจสอบความพร้อมของหน้าจอหลักเมื่อเปิดใช้งาน}
เมื่อเปิดโปรแกรม PathologyApp ขึ้นมาครั้งแรก ระบบจะทำการโหลดแบบจำลองภาษาและเครื่องมือสกัดข้อมูลออฟไลน์เข้าสู่หน่วยความจำหลักโดยอัตโนมัติ โดยผู้ใช้งานจะพบกับหน้าจอหลักที่พร้อมรองรับการบันทึกเสียงบรรยายผลการตรวจชิ้นเนื้อ 15 ฟิลด์ตามมาตรฐาน CAP

\begin{figure}[htbp]
    \centering
    \includegraphics[width=0.85\textwidth]{คู่มือการติดตั้ง/Screenshot 2026-10-22 163428.png}
    \caption{แสดงหน้าจอหลัก (User Interface) ของแอปพลิเคชัน PathologyApp พร้อมใช้งานแบบออฟไลน์}
    \label{fig:fig6}
\end{figure}

\appendiceschapter{คู่มือการใช้งานระบบแอปพลิเคชัน}
\label{app:usermanual}

ภาคผนวกนี้แสดงขั้นตอนและคู่มือการใช้งานระบบแอปพลิเคชัน ระบบแปลงเสียงเป็นข้อความ สำหรับสนับสนุนการบันทึกข้อมูลการตรวจสิ่งส่งตรวจทางศัลยพยาธิวิทยาด้วยตาเปล่า (Gross Examination) ครอบคลุมตั้งแต่การเปิดระบบจากหน้าจอเดสก์ท็อป การเข้าสู่ระบบ การเลือกโหมดการทำงาน การบันทึกเสียงสดหรืออัปโหลดไฟล์เสียง การถอดความและสกัดข้อมูลลงแบบฟอร์มอัตโนมัติ การตรวจทานแก้ไขร่วมกับระบบธงเตือนความปลอดภัย ไปจนถึงการส่งออกเอกสารรายงานผลทางการแพทย์ในรูปแบบ PDF

\section*{ขั้นตอนที่ 1: การเริ่มต้นระบบจากเดสก์ท็อป (System Startup)}
ผู้ใช้งานเริ่มต้นระบบโดยดับเบิลคลิกไอคอนทางลัด \textbf{Pathology Voice Assistant} บนหน้าจอเดสก์ท็อปของเครื่องคอมพิวเตอร์ประจำสถานีตรวจชิ้นเนื้อ (Grossing Station) ดังแสดงในภาพที่ \ref{fig:man_shortcut} ระบบจะเริ่มการทำงานของเซิร์ฟเวอร์แบบออฟไลน์และเปิดเบราว์เซอร์เข้าสู่หน้าจอการทำงานโดยอัตโนมัติ

\begin{figure}[H]
\centering
\includegraphics[width=0.75\textwidth]{images/manual/manual_1_shortcut.png}
\caption{การเริ่มต้นระบบผ่านไอคอนทางลัดบนหน้าจอเดสก์ท็อป (Desktop Shortcut)}
\label{fig:man_shortcut}
\end{figure}

\section*{ขั้นตอนที่ 2: การเข้าสู่ระบบและการเลือกระดับความซับซ้อน (Sign In)}
เมื่อเข้าสู่ระบบ จะปรากฏหน้าต่างดังภาพที่ \ref{fig:man_login} ผู้ใช้งานกรอกชื่อผู้ตรวจ (เช่น \texttt{user-1}) และเลือกระดับความซับซ้อนของสิ่งส่งตรวจ (เช่น \textit{Level 1 - กรณีทั่วไป / ไม่ซับซ้อน (Routine / Simple)}) จากนั้นคลิกปุ่ม \textbf{เข้าสู่ระบบ / เริ่มบันทึกเคส} เพื่อเข้าสู่หน้าจอหลัก

\begin{figure}[H]
\centering
\includegraphics[width=0.75\textwidth]{images/manual/manual_2_login.png}
\caption{หน้าจอเข้าสู่ระบบและการระบุระดับความซับซ้อนของกรณีศึกษา}
\label{fig:man_login}
\end{figure}

\section*{ขั้นตอนที่ 3: หน้าจอแดชบอร์ดหลักและการเลือกแบบฟอร์ม (Dashboard)}
หน้าจอแดชบอร์ดหลัก (ภาพที่ \ref{fig:man_dashboard}) แสดงตัวเลือกการทำงาน 2 รูปแบบหลัก ได้แก่:
\begin{enumerate}[label=(\arabic*), leftmargin=*, align=left]
    \item \textbf{การเริ่มกรณีศึกษาใหม่ด้วยการบันทึกเสียงสด (START A CASE):} คลิกที่ปุ่ม \textbf{New case $\rightarrow$} ในการ์ด \textit{Breast gross} เพื่อเปิดแบบฟอร์มการตรวจชิ้นเนื้อเต้านมมาตรฐานสำหรับการบรรยายเสียงแบบเรียลไทม์
    \item \textbf{การเริ่มกรณีศึกษาจากไฟล์เสียงที่บันทึกไว้ล่วงหน้า (START FROM A RECORDING):} ลากและวางไฟล์เสียง หรือคลิกปุ่ม \textbf{Upload and open the case} เพื่อนำเข้าไฟล์เสียงบันทึกภายนอก
\end{enumerate}
นอกจากนี้ ที่ส่วนล่างของหน้าจอจะแสดงแถบ \textit{Clinical AI Data Flywheel} สำหรับติดตามสถิติการรวบรวมคลังข้อมูลเสียงจริงที่ผ่านการตัดข้อมูลระบุตัวบุคคลออกทั้งหมด (De-identified 100\%) สำหรับใช้ในการวิจัยและพัฒนาต่อยอดในอนาคต

\begin{figure}[H]
\centering
\includegraphics[width=0.88\textwidth]{images/manual/manual_3_dashboard.png}
\caption{หน้าจอแดชบอร์ดหลักแสดงตัวเลือกการสร้างกรณีศึกษาใหม่และการนำเข้าไฟล์เสียง}
\label{fig:man_dashboard}
\end{figure}

\section*{ขั้นตอนที่ 4: การเปิดกรณีศึกษาผ่านการอัปโหลดไฟล์เสียงบันทึก (Audio Upload)}
ในกรณีที่เลือกการนำเข้าไฟล์เสียงจากเครื่องบันทึกเสียงพกพา ระบบจะแสดงหน้าต่างเลือกไฟล์ (File Picker Dialog) ดังภาพที่ \ref{fig:man_upload} ผู้ใช้สามารถเลือกไฟล์เสียงบันทึกการตรวจชิ้นเนื้อ (รองรับรูปแบบ \texttt{.mp3}, \texttt{.m4a}, \texttt{.wav}) จากนั้นคลิกปุ่ม \textbf{เปิด} เพื่อส่งไฟล์เข้าสู่กระบวนการประมวลผลทันที

\begin{figure}[H]
\centering
\includegraphics[width=0.88\textwidth]{images/manual/manual_5_upload.png}
\caption{หน้าต่างการเลือกและอัปโหลดไฟล์เสียงบันทึกการตรวจชิ้นเนื้อเข้าสู่ระบบ}
\label{fig:man_upload}
\end{figure}

\section*{ขั้นตอนที่ 5: หน้าจอแบบฟอร์มเปล่าและการเตรียมบันทึกเสียงสด (Blank Form View)}
ในกรณีที่เปิดกรณีศึกษาใหม่สำหรับการบันทึกเสียงสด หน้าจอจะแสดงแบบฟอร์มการตรวจชิ้นเนื้อเต้านมมาตรฐาน 15 ฟิลด์ (\textit{Breast gross 1}) ทางฝั่งซ้าย และแผงควบคุมการบันทึกเสียงทางฝั่งขวา ดังภาพที่ \ref{fig:man_blank} โดยมีสถานะไมโครโฟนแสดงความพร้อมทำงาน (\texttt{Ready}) เมื่อพยาธิแพทย์พร้อมบรรยาย ให้คลิกปุ่ม \textbf{Record} เพื่อเริ่มบันทึกเสียงพูด

\begin{figure}[H]
\centering
\includegraphics[width=0.88\textwidth]{images/manual/manual_4_blank_form.png}
\caption{หน้าจอแบบฟอร์มเปล่าสำหรับการตรวจชิ้นเนื้อเต้านมและแผงควบคุมการบันทึกเสียงสด}
\label{fig:man_blank}
\end{figure}

\section*{ขั้นตอนที่ 6: การประมวลผลถอดความ การสกัดข้อมูล และระบบธงเตือนความปลอดภัย}
เมื่อสิ้นสุดการบรรยาย (คลิกปุ่ม \textbf{Stop}) ระบบจะประเมินค่าอัตราส่วนสัญญาณต่อเสียงรบกวนแบบไม่ต้องใช้สัญญาณอ้างอิง ($\widehat{\mathrm{SNR}}$) เพื่อคัดกรองสัญญาณผ่านกลไก \texttt{D8\_snr\_gated} เฉพาะเมื่อตรวจพบเสียงรบกวน ($\widehat{\mathrm{SNR}} < 13.55\text{ dB}$) และทำการถอดความด้วยแบบจำลอง \textbf{PathoWhisper Small LoRA (INT8)} ร่วมกับข้อความนำบริบทที่ระบุคู่คำพ้องเสียง ($P_4$) ข้อความถอดความเบื้องต้น (Transcript) จะแสดงในกล่องข้อความฝั่งขวา ดังภาพที่ \ref{fig:man_transcribe} จากนั้นตัวปรับข้อความและตัวสกัดข้อมูลรุ่นที่ 3 (\texttt{extractor\_v3.py}) จะบรรจุค่ามิติ ขนาดก้อนเนื้องอก ตำแหน่ง และตัวเลือกในแบบฟอร์มมาตรฐานทั้ง 15 ฟิลด์ทางฝั่งซ้ายโดยอัตโนมัติ

\begin{figure}[H]
\centering
\includegraphics[width=0.88\textwidth]{images/manual/manual_6_transcription_fill.png}
\caption{การประมวลผลถอดความเสียงและการจัดสรรข้อมูลลงแบบฟอร์มอัตโนมัติ พร้อมการแสดงแถบธงเตือนความปลอดภัย}
\label{fig:man_transcribe}
\end{figure}

\noindent \textbf{การทำงานของระบบความปลอดภัยหลายสัญญาณและการตรวจทานโดยพยาธิแพทย์:}
\begin{itemize}
    \item \textbf{แถบแจ้งเตือนความปลอดภัยทางคลินิกหลายสัญญาณ (Multi-Signal Clinical Warning Banners):} ระบบผสานค่าความมั่นใจระดับประโยคและระดับคำของตัวถอดรหัสเข้ากับกฎตรวจสอบความสอดคล้องทางคลินิก และแสดงแถบแจ้งเตือนสีส้มที่ด้านบนของแบบฟอร์ม (ภาพที่ \ref{fig:man_transcribe}) เช่น แจ้งเตือนกรณีเลือกหัตถการผ่าตัดแบบ MRM แต่ยังไม่ได้ระบุข้อมูลต่อมน้ำเหลืองรักแร้ (\textit{Clinical Procedure Check}), แจ้งเตือนกรณีพบรอยโรคแต่ขาดข้อมูลขนาดก้อนเนื้องอกหรือระยะห่างขอบตัดด้านลึก, หรือแจ้งเตือนกรณีไม่พบเลขที่สิ่งส่งตรวจ (\textit{Missing Surgical Number})
    \item \textbf{การตรวจทานและแก้ไขโดยแพทย์ (Physician-in-the-Loop Verification):} พยาธิแพทย์ผู้ตรวจสามารถคลิกแก้ไขตัวเลข มิติ หรือข้อความในช่องข้อมูลบนแบบฟอร์มได้โดยตรง หรือแก้ไขข้อความถอดความแล้วคลิกปุ่ม \textbf{Fill again} เพื่อสั่งประมวลผลการสกัดข้อมูลซ้ำ
\end{itemize}

\section*{ขั้นตอนที่ 7: การบันทึกข้อมูลและการส่งออกเอกสารรายงานผล PDF}
เมื่อพยาธิแพทย์ผู้ตรวจตรวจสอบและยืนยันความถูกต้องของข้อมูลครบถ้วนแล้ว ให้คลิกปุ่มบันทึกข้อมูลเพื่อจัดเก็บลงฐานข้อมูลภายในเครื่อง (PostgreSQL ร่วมกับฐานข้อมูลสำรอง SQLite) จากนั้นคลิกปุ่ม \textbf{Export PDF} เพื่อสร้างเอกสารรายงานผลทางการแพทย์มาตรฐาน ดังแสดงในภาพที่ \ref{fig:man_pdf} เอกสารที่สร้างขึ้นจะมีโครงสร้างตารางมาตรฐานครบทั้ง 15 ส่วน พร้อมเลขที่สิ่งส่งตรวจและรหัสอ้างอิงประจำกรณีศึกษาสำหรับใช้ตรวจสอบย้อนกลับ และพร้อมสำหรับการพิมพ์แนบกับสิ่งส่งตรวจหรือนำเข้าสู่ระบบสารสนเทศของโรงพยาบาล

\begin{figure}[H]
\centering
\includegraphics[width=0.85\textwidth]{images/manual/manual_7_export_pdf.png}
\caption{เอกสารรายงานผลการตรวจชิ้นเนื้อเต้านมมาตรฐานในรูปแบบ PDF สำหรับพิมพ์หรือนำเข้าสู่ระบบสารสนเทศของโรงพยาบาล}
\label{fig:man_pdf}
\end{figure}

\noindent \textbf{หมายเหตุเกี่ยวกับส่วนต่อประสานเพิ่มเติม:}
บนหน้าจอส่วนต่อประสานผู้ใช้มีปุ่มควบคุมฟังก์ชันกล้องถ่ายภาพ (\texttt{Photo / Front Camera}) และแถบสถานะการตรวจจับท่าทางมือ (\texttt{Gesture Control}) ซึ่งเป็นส่วนประกอบของระบบต้นแบบที่ออกแบบเตรียมไว้สำหรับรองรับการเชื่อมต่อกล้องถ่ายภาพสิ่งส่งตรวจและการสั่งงานแบบไร้การสัมผัส (Touchless Hands-Free Control) ในอนาคต ทั้งนี้ ฟังก์ชันเสริมดังกล่าวอยู่นอกเหนือขอบเขตการประเมินประสิทธิภาพของระบบรู้จำเสียงพูดและการสกัดข้อมูลในรายงานฉบับนี้

\appendiceschapter{การเปิดเผยการใช้งานปัญญาประดิษฐ์และเครื่องมือช่วยพัฒนา}
\label{app:aideclaration}

เพื่อความโปร่งใสทางวิชาการและเป็นไปตามแนวปฏิบัติการใช้งานปัญญาประดิษฐ์เชิงสร้างสรรค์ (Generative AI) ในการทำโครงงานวิจัยและวิทยานิพนธ์ ผู้จัดทำได้เปิดเผยรายละเอียดการประยุกต์ใช้โมเดลภาษาขนาดใหญ่ (Large Language Models: LLM) และตัวแทนปัญญาประดิษฐ์ (AI Agents) ในขั้นตอนการวิจัยและพัฒนาระบบ PathoWhisper โดยการตัดสินใจทางสถาปัตยกรรม การออกแบบการทดลอง การทวนสอบความถูกต้องของซอร์สโค้ด และการรับผิดชอบต่อความถูกต้องของเนื้อหาทั้งหมดในเล่มวิทยานิพนธ์ฉบับนี้ เป็นการดำเนินการและกำกับดูแลโดยผู้จัดทำโครงงาน (Human-in-the-Loop) ทั้งสิ้น ดังแสดงรายการคำสั่งและบทบาทการควบคุมในแต่ละภารกิจต่อไปนี้:

\section*{รายการภารกิจ คำสั่ง Prompt และการควบคุมโดยมนุษย์}

\begin{enumerate}[label=\textbf{\arabic*.}, leftmargin=*, align=left]

\item \textbf{การออกแบบและสังเคราะห์ชุดข้อมูลทางการแพทย์ (Synthetic Clinical Dataset Generation)}
\begin{itemize}[leftmargin=1.5em]
    \item \textbf{เครื่องมือที่ใช้:} Anthropic Claude 3.5 Sonnet และ OpenAI GPT-4o
    \item \textbf{ตัวอย่าง Prompt / คำสั่ง:} ``เสนอไอเดียและช่วยร่างโครงสร้างประโยคการบรรยายผลตรวจชิ้นเนื้อมะเร็งเต้านมด้วยตาเปล่า (Gross Examination) ตามเกณฑ์มาตรฐาน CAP ให้ครอบคลุม 10 หมวดหมู่ทางคลินิก โดยให้แยกคลังสำนวนประโยคสำหรับชุดฝึกฝนและชุดทดสอบอย่างเด็ดขาดเพื่อป้องกัน Data Leakage''
    \item \textbf{การควบคุมและปรับปรุงโดยมนุษย์:} คัดเลือกและตรวจสอบความถูกต้องของคำศัพท์ทางการแพทย์ร่วมกับข้อกำหนดของ College of American Pathologists (CAP Protocol) และคัดกรองความสมจริงของมิติขนาดสิ่งส่งตรวจและก้อนมะเร็ง
\end{itemize}

\item \textbf{การพัฒนาอัลกอริทึมคัดกรองสัญญาณรบกวนแบบปรับตัว (Blind SNR-Gated Adaptive Denoising)}
\begin{itemize}[leftmargin=1.5em]
    \item \textbf{เครื่องมือที่ใช้:} Google DeepMind Antigravity (AI Coding Assistant)
    \item \textbf{ตัวอย่าง Prompt / คำสั่ง:} ``ช่วยเขียนฟังก์ชันคำนวณ Blind Frame-Energy Percentile SNR ($\widehat{\mathrm{SNR}} = P_{85} - P_{15}$) จากสัญญาณเสียงในภาษา Python และเขียนตรรกะ FFmpeg คำสั่ง afftdn สำหรับเปิดปิดตัวกรองความถี่เฉพาะเมื่อค่า SNR ต่ำกว่าเกณฑ์ที่กำหนดเพื่อไม่ให้กระทบเสียงพูดสะอาด''
    \item \textbf{การควบคุมและปรับปรุงโดยมนุษย์:} ตรวจสอบสูตรคณิตศาสตร์ ทำการทดลองสอบเทียบเกณฑ์ขีดแบ่ง ($\tau^* = 13.55\text{ dB}$) ด้วย Youden's Index บนชุดตรวจสอบจริง และทดลองค้นหาแบบกริดเพื่อเลือกพารามิเตอร์ของตัวกรอง ($\mathrm{nr} = 6\text{ dB}, \mathrm{nf} = -30\text{ dB}$) ที่ไม่ทำลายเสียงพยัญชนะสั้น
\end{itemize}

\item \textbf{การออกแบบข้อความนำบริบทแบบเปรียบเทียบเสียงพ้อง (Phonetic-Contrastive Prompt Engineering)}
\begin{itemize}[leftmargin=1.5em]
    \item \textbf{เครื่องมือที่ใช้:} Anthropic Claude 3.5 Sonnet
    \item \textbf{ตัวอย่าง Prompt / คำสั่ง:} ``ช่วยวิเคราะห์คำคู่พ้องเสียง (Homophones) ทางการแพทย์ที่โมเดล Whisper มักถอดความสับสนในงานตรวจชิ้นเนื้อ เช่น right/write, discrete/discreet, deep surgical margin/lymph margin และช่วยร่างโครงสร้าง Prompt $P_4$ ให้อยู่ในงบประมาณ 224 โทเคน''
    \item \textbf{การควบคุมและปรับปรุงโดยมนุษย์:} ตรวจสอบความถูกต้องของคำศัพท์แพทย์ ดำเนินการทดลองเปรียบเทียบข้อความนำบริบท ($P_0 \dots P_4$) แบบเป็นระบบ (Prompt Ablation Study) และวิเคราะห์การแยกองค์ประกอบความผิดพลาดเชิงตั้งฉากระหว่าง C-TER และ N-DER
\end{itemize}

\item \textbf{การปรับแต่งแบบจำลอง LoRA และการลดทอนความละเอียด (Whisper Small LoRA Fine-Tuning \& INT8 Quantization)}
\begin{itemize}[leftmargin=1.5em]
    \item \textbf{เครื่องมือที่ใช้:} Google DeepMind Antigravity (AI Coding Assistant)
    \item \textbf{ตัวอย่าง Prompt / คำสั่ง:} ``ช่วยเขียนโค้ด PyTorch ร่วมกับไลบรารี HuggingFace PEFT เพื่อทำ Parameter-Efficient Fine-Tuning (LoRA: $r=8, \alpha=32$) บนเลเยอร์ Decoder Attention ของแบบจำลอง Whisper Small พร้อมสคริปต์ควบรวมน้ำหนักและแปลงเป็น CTranslate2 INT8 เพื่อรันบน CPU''
    \item \textbf{การควบคุมและปรับปรุงโดยมนุษย์:} ตรวจสอบโครงสร้างเลเยอร์เป้าหมาย ($48$ โมดูล, $589,824$ พารามิเตอร์ หรือ $0.2434\%$), มอนิเตอร์กราฟ Loss ระหว่างการฝึกฝนเพื่อควบคุมไม่ให้เกิด Overfitting ที่ 3 Epochs และทดสอบวัดอัตราความเร็ว RTF บนเครื่องคอมพิวเตอร์ส่วนบุคคล
\end{itemize}

\item \textbf{การพัฒนาระบบสกัดข้อมูล 15 ฟิลด์ และตัวจัดการการพูดแก้ไขตัวเอง (Deterministic NLP Extractor \& Mid-Sentence Self-Correction)}
\begin{itemize}[leftmargin=1.5em]
    \item \textbf{เครื่องมือที่ใช้:} Google DeepMind Antigravity / Anthropic Claude 3.5 Sonnet
    \item \textbf{ตัวอย่าง Prompt / คำสั่ง:} ``ช่วยออกแบบ Regular Expressions และโครงสร้าง Contextual Pattern Matching สำหรับสกัดข้อมูล 15 ฟิลด์ CAP โดยเพิ่มกลไกตรวจจับการพูดแก้ไขตัวเองกลางประโยค เช่น คำว่า "sorry", "แก้เป็น", "wait" และการแยกแยะคำว่า "weight/grams" ออกจากคำว่า "wait"''
    \item \textbf{การควบคุมและปรับปรุงโดยมนุษย์:} ออกแบบและรันชุดทดสอบระดับหน่วย (Unit Test: TC-01 ถึง TC-16) ทวนสอบความถูกต้องของการสกัดข้อมูลจนได้ Macro-$F_1 = 100.0\%$ บนข้อความเฉลย และตรวจสอบกรณีประโยคปฏิเสธทางคลินิก (Negation Handling)
\end{itemize}

\item \textbf{การพัฒนาระบบธงเตือนความปลอดภัยทางคลินิกหลายสัญญาณ (Multi-Signal Clinical Safety Flagging)}
\begin{itemize}[leftmargin=1.5em]
    \item \textbf{เครื่องมือที่ใช้:} Google DeepMind Antigravity
    \item \textbf{ตัวอย่าง Prompt / คำสั่ง:} ``ช่วยเขียนฟังก์ชันคำนวณคะแนนความเสี่ยงหลายสัญญาณ โดยผสานค่าความน่าจะเป็นทางเสียง (Average Logprob, Minimum Word Probability) เข้ากับกฎความสอดคล้องทางคลินิก 15 ฟิลด์ และเขียนอัลกอริทึม Selective Prediction สำหรับวิเคราะห์ Risk-Coverage''
    \item \textbf{การควบคุมและปรับปรุงโดยมนุษย์:} กำหนดเกณฑ์ขีดแบ่งความเสี่ยง ($\tau_{\text{risk}} = 0.50$) ทำการคำนวณ Risk-Coverage Trade-off และ Selective Classification บนข้อมูลทดสอบ 300 กรณี และวิเคราะห์สาเหตุของข้อผิดพลาดหลุดรอดแบบเงียบ (Silent CER)
\end{itemize}

\item \textbf{การพัฒนาส่วนต่อประสานแอปพลิเคชันและการทำงานแบบออฟไลน์ (Web Application, Hands-Free UI \& Offline Architecture)}
\begin{itemize}[leftmargin=1.5em]
    \item \textbf{เครื่องมือที่ใช้:} Google DeepMind Antigravity (AI Coding Assistant)
    \item \textbf{ตัวอย่าง Prompt / คำสั่ง:} ``ช่วยพัฒนาระบบแอปพลิเคชันด้วย Flask สำหรับสถานีตรวจชิ้นเนื้อ รองรับการรับเสียงไมโครโฟนแบบ Hands-Free, การแสดง Audio Waveform แบบเคลื่อนไหว, การสลับใช้งานโมเดล CTranslate2 ออฟไลน์ ''
    \item \textbf{การควบคุมและปรับปรุงโดยมนุษย์:} ตรวจสอบความปลอดภัยของซอร์สโค้ด ตรวจสอบการจัดการทรัพยากรหน่วยความจำของ CPU ทดสอบฟังก์ชันการทำงานในสภาวะตัดการเชื่อมต่ออินเทอร์เน็ต (Isolated Offline Testing) และปรับแต่ง UI ให้เหมาะสมกับการสวมถุงมือปฏิบัติงานจริง
\end{itemize}

\item \textbf{การพัฒนาระบบสร้างเอกสารรายงานทางการแพทย์อัตโนมัติ (Automated Pathology PDF Report Generation)}
\begin{itemize}[leftmargin=1.5em]
    \item \textbf{เครื่องมือที่ใช้:} Google DeepMind Antigravity
    \item \textbf{ตัวอย่าง Prompt / คำสั่ง:} ``ช่วยเขียนฟังก์ชันด้วย PyMuPDF (fitz) เพื่อประมวลผลวาดเครื่องหมายกากบาท/ถูก และพิมพ์ข้อความสกัด 15 ฟิลด์ลงในแม่แบบเอกสาร PDF Breast Gross Examination ให้ตรงพิกัดและขนาดช่องตารางของโรงพยาบาล''
    \item \textbf{การควบคุมและปรับปรุงโดยมนุษย์:} วัดพิกัด Bounding Box บนแบบฟอร์มจริง ปรับแต่งรูปแบบการตัดคำและการจัดวางฟอนต์ภาษาไทยและอังกฤษให้คมชัด สวยงาม และตรงตามมาตรฐานเวชระเบียน
\end{itemize}

\item \textbf{การหาจุดบกพร่องของโค้ดและการทดสอบระดับหน่วย (Code Debugging \& Unit Testing)}
\begin{itemize}[leftmargin=1.5em]
    \item \textbf{เครื่องมือที่ใช้:} Google DeepMind Antigravity (AI Coding Assistant)
    \item \textbf{ตัวอย่าง Prompt / คำสั่ง:} ``ช่วยสแกนหาจุดบกพร่องในสคริปต์ถอดรหัสเสียงยาว fix\_longform\_lora\_eval.py ที่พบอาการวนซ้ำคำ (Repetition Loop) ในกรณีเสียงเกิน 27.5 วินาที และช่วยร่างโค้ดเพิ่ม repetition\_penalty=1.15 และ vad\_filter=True''
    \item \textbf{การควบคุมและปรับปรุงโดยมนุษย์:} ทำการ Debugging ตรวจสอบ Log และรันชุดทดสอบความถูกต้องจนขจัดปัญหาการถอดความค้างในเคสยาวได้สำเร็จ 100\%
\end{itemize}
\end{enumerate}

\clearpage
\addcontentsline{toc}{chapter}{บรรณานุกรม}
\makeatletter
\let\ps@plain\ps@empty
\makeatother
\bibliographystyle{ieeetr}
\bibliography{references}

\end{document}